\documentclass[10pt,a4paper]{amsart}
\usepackage[margin=19mm]{geometry}
\usepackage{amsmath,amssymb,mathtools}
\usepackage[scr=rsfs,cal=euler]{mathalfa}
\usepackage{xfrac}
\usepackage[unicode,hidelinks,bookmarks=false]{hyperref}
\newcommand{\ie}{i.e.\ }
\newcommand{\cf}{cf.\ }
\newcommand{\resp}{respectively\ }
\newcommand{\Subsection}{\subsection}
\providecommand{\nobreakdash}{\nobreak\hbox{-}}
\setlength{\emergencystretch}{2em}
\multlinegap0pt
\usepackage{bm}
\usepackage{bbm}
\usepackage{dsfont}
\usepackage{xspace}
\usepackage{cancel}
\usepackage{mathtools}
\usepackage{amsthm}
\makeatletter
\@ifundefined{lemma}{\newtheorem{lemma}{Lemma}}{}
\makeatother
\newcommand{\indep}{\rotatebox[origin=c]{90}{$\models$}}
\title[On the Benefits of Accelerated Optimization in Robust and Pr]{On the Benefits of Accelerated Optimization in Robust and Private Estimation}
\author{Laurentiu Andrei Marchis, Po-Ling Loh}
\date{Source: arXiv:2506.03044v1 (CC BY 4.0)}
\begin{document}
\maketitle

\noindent\emph{Source and licence.} On the Benefits of Accelerated Optimization in Robust and Private Estimation. Version arXiv:2506.03044v1 (CC BY 4.0), distributed under the Creative Commons Attribution 4.0 International licence, as recorded in the arXiv licence metadata for this exact version and shown on the abstract page https://arxiv.org/abs/2506.03044v1. Full rights notice: licenses/arxiv-2506.03044-CC-BY-4.0.txt.\par\smallskip

\noindent\textit{Editorial scope.} Counted unit: the Lemma whose source label is \texttt{lemma:pHu\_prelim} in the article (source0.tex lines 1677--1698, source label \texttt{lemma:pHu\_prelim}), delivered together with the article's own complete original proof of it (source0.tex lines 1700--1765, one \texttt{proof} environment of 66 lines). No mathematics is written, repaired or re-derived: statement and proof are the article's own text line for line. Required context retained uncounted: none. Every definition, display and label of those lines is retained unchanged. Omitted from this package and not redistributed: 1--1676, 1699--1699, 1766--3806. Nothing inside a retained excerpt is omitted or altered: every retained body is delivered exactly as the article's lines are written, so no omission receipt of any kind accompanies this package. Citations: the retained excerpts carry no citation command at all, so no bibliography entry of the article is transcribed and no reference is redistributed. Reference adaptation: none was needed and none was made; every reference inside the delivered unit resolves inside it, so no unresolved reference or citation is produced. Printed slips kept literal: no printed word, symbol, hypothesis or number of the counted statement or of its proof was altered. Layout and class: the article's own class is \texttt{article} with packages including algorithm, algpseudocode, adjustbox, indentfirst, amsthm, inputenc, mathtools, amsfonts, array, booktabs, multirow, makecell, float, etoolbox; the wrapper is the lane's pinned \texttt{amsart} editorial wrapper at 10pt on A4, which restates the theorem-like environments the retained text needs and adds the article's own macro definitions that the retained text needs (\texttt{\textbackslash indep}), with extra wrapper packages (bm, bbm, dsfont, xspace, cancel, mathtools). The delivered numbering is the unit's own. Assets: no retained span contains a figure environment, a graphics-inclusion command, a PDF-page inclusion, a table or a TikZ picture; no image file is copied into this package, the package declares no asset dependency and no image file is redistributed with it. Licence: version arXiv:2506.03044v1 of this article is distributed under the Creative Commons Attribution 4.0 International licence, as recorded in the arXiv licence metadata for this exact version and shown on the abstract page https://arxiv.org/abs/2506.03044v1. A full rights notice is delivered at licenses/arxiv-2506.03044-CC-BY-4.0.txt. Characters: every retained excerpt is pure ASCII, so the delivered engine, XeLaTeX with its default Latin Modern text font, reports no missing character.
\begin{lemma}
\label{lemma:pHu_prelim}
Let $L_x, C''_1, q > 0$ and $C'_2 \geq C'_1 > 0$. On the domain $||x||_2 \leq L_x \ and \ y \in \mathbb{R}$, define the loss
\begin{align}
\label{eq:ps_Hub_loss}
\mathcal{L}(\theta, (x, y)) = \rho_q(y - x^T\theta), \quad \forall \theta \in \mathbb{R}^p.
\end{align}
Then:
\begin{enumerate}
\item $\mathcal{L}$ is $qL_x$-Lipschitz in $\theta$.

\item Consider the linear regression model $y = x^T\theta^* + w$, with $\mathbb{E}[x] = 0$, $\mathbb{E}[w] = 0$, $\Sigma = \mathbb{E}[xx^T]$, and $x \indep w$. Assume $\lambda_{\max}(\Sigma) \leq \frac{C'_2}{p}$. Then the corresponding risk $\mathcal{R}$ to \eqref{eq:ps_Hub_loss} is $\frac{C'_2}{p}$-smooth over $\mathbb{R}^p$. 

\item Additionally, let $\mathcal{C}$ be a convex set such that $\theta^* \in \mathcal{C}$ and $||\mathcal{C}||_2 \leq C''_1\sqrt{p}$. Assume $\frac{C'_1}{p} \leq \lambda_{\min}(\Sigma)$ and $x$ has bounded $4^{\text{th}}$ moments, i.e., there exists $\widetilde{C}_4 > 0$ such that $\mathbb{E}\left[(x^Tv)^4\right] \leq \widetilde{C}_4\mathbb{E}\left[(x^Tv)^2\right]^2$, for any $||v||_2 = 1$. Then the risk is 
\begin{align*}
\frac{q^3(C'_1)^4}{4p\left((C'_1)^2q^2 + 8(C''_1)^2(C'_2)^3\widetilde{C}_4 + 2(C'_1)^2\sigma_2^2\right)^{3/2}}\mbox{-strongly convex}
\end{align*}
over $\mathcal{C}$. 

\item $\nabla\mathcal{R}(\theta^*) = 0$, so $\theta_{*} = \theta^*$ is the minimizer of $\mathcal{R}$ over $\mathcal{C}$.
\end{enumerate}
\end{lemma}

\begin{proof}
We first prove (1).
%Also, for now, $(x, y) \in \mathbb{B}_2(L_x) \times \mathbb{R}$, not necessarily from the linear regression model: $y = x^T\theta^* + w$, with $\mathbb{E}[x] = 0$, $\mathbb{E}[w] = 0$, $\Sigma = \mathbb{E}[xx^T]$, and $x \indep w$.
Note that
\begin{align*}
\nabla \mathcal{L}(\theta, (x, y)) = -\psi_q(y - x^T\theta)x, \ \forall \theta \in \mathbb{R}^p. 
\end{align*}
Hence, we clearly have $||\nabla \mathcal{L}(\theta, (x, y))||_2 \leq q L_x$ on the domain.
%So, $\mathcal{L}$ is $qL_x$-Lipschitz in $\theta \in \mathbb{R}^p$, for all $||x||_2 \leq L_x$ and $y \in \mathbb{R}$. This proves the first part, i.e., the Lipschitz property.

%Now, consider the linear regression model: $y = x^T\theta^* + w$, with $\mathbb{E}[x] = 0$, $\mathbb{E}[w] = 0$, $\Sigma = \mathbb{E}[xx^T]$, and $x \indep w$.
For (2), note that $\nabla\mathcal{R}(\theta) = -\mathbb{E}[\psi_q(y - x^T\theta)x]$, for all $\theta \in \mathbb{R}^p$. Since $\psi_q$ is bounded and differentiable, we can swap expectations and derivatives by the Dominated Convergence Theorem to obtain 
\begin{align*}
\nabla^2\mathcal{R}(\theta) = \mathbb{E}[\psi'_q(y - x^T\theta)xx^T].
\end{align*}
Take $\theta \in \mathbb{R}^p$. Note that $0 < \psi'_q(t) \leq 1$, for $t \in \mathbb{R}$. Hence, we have
\begin{align*}
\nabla^2\mathcal{R}(\theta) = \mathbb{E}[\psi'_q(y - x^T\theta)xx^T] \preceq \Sigma \preceq \lambda_{\max}(\Sigma)I_p \preceq \frac{C'_2}{p}I_p.
\end{align*}

%Now, additionally, take $\theta \in \mathcal{C}$, with $||\mathcal{C}||_2 \leq C''_1\sqrt{p}$, and assume $x$ has bounded $4^{th}$ moments as per Definition \ref{def:bd_moments}, and that $\frac{C'_1}{p} \leq \lambda_{\min}(\Sigma)$.
For (3), let $a := \theta^* - \theta$. By Markov's inequality, since $x \indep w$, we have
\begin{align*}
\nabla^2\mathcal{R}(\theta) &= \mathbb{E}\left[\frac{1}{\left(1 + \left(\frac{x^Ta + w}{q}\right)^2\right)^{3/2}}xx^T\right] \\
%
& \succeq \mathbb{E}\left[\frac{1}{\left(1 + \left(\frac{|x^Ta| + |w|}{q}\right)^2\right)^{3/2}}xx^T\bigg\rvert|w| < 2\mathbb{E}[|w|]\right]\mathbb{P}(|w| < 2\mathbb{E}[|w|])\\
&\succeq \frac{1}{2}\mathbb{E}\left[\frac{1}{\left(1 + \left(\frac{|x^Ta| + \mathbb{E}[|w|]}{q}\right)^2\right)^{3/2}}xx^T\right]\\ 
&= \frac{1}{2}\mathbb{E}\left[\frac{q^3}{\left(q^2 + \left(|x^Ta| + \mathbb{E}[|w|]\right)^2\right)^{3/2}}xx^T\right].
\end{align*}
Let $C'_3 = \frac{2C'_2\sqrt{\widetilde{C}_4}}{C'_1}$ and $A = \left\{|x^Ta| < C'_3||\mathcal{C}||_2\sqrt{\lambda_{\max}(\Sigma)}\right\}$. Again using Markov's inequality, we obtain
\begin{align*}
\mathbb{P}(A^c) \leq \frac{\mathbb{E}\left[(x^Ta)^2\right]}{(C'_3)^2||\mathcal{C}||_2^2\lambda_{\max}(\Sigma)} = \frac{a^T\Sigma a}{(C'_3)^2||\mathcal{C}||_2^2\lambda_{\max}(\Sigma)} \leq \frac{||a||_2^2\lambda_{\max}(\Sigma)}{(C'_3)^2||\mathcal{C}||_2^2\lambda_{\max}(\Sigma)} \leq \frac{1}{(C'_3)^2},
\end{align*}
where $A^c$ denotes the complement of $A$. Hence, we have
\begin{align*}
\nabla^2\mathcal{R}(\theta) &\succeq \frac{1}{2}\mathbb{E}\left[\frac{q^3}{\left(q^2 + \left(|x^Ta| + \mathbb{E}[|w|]\right)^2\right)^{3/2}}xx^T\mathbbm{1}_A\right]\\
&\succeq \frac{q^3}{2\left(q^2 + \left(C'_3||\mathcal{C}||_2\sqrt{\lambda_{\max}(\Sigma)} + \mathbb{E}[|w|]\right)^2\right)^{3/2}}\left(\mathbb{E}[xx^T] - \mathbb{E}[xx^T\mathbbm{1}_{A^c}]\right).
\end{align*}
Take $||v||_2 = 1$ arbitrary. We have by Cauchy-Schwarz that
\begin{align*}
v^T\left(\mathbb{E}[xx^T] - \mathbb{E}[xx^T\mathbbm{1}_{A^c}]\right)v &\geq \lambda_{\min}(\Sigma) - \mathbb{E}\left[(x^Tv)^2\mathbbm{1}_{A^c}\right]\\
&\geq \lambda_{\min}(\Sigma) - \sqrt{\mathbb{E}\left[(x^Tv)^4\right]\mathbb{P}(A^c)}\\
&\geq \lambda_{\min}(\Sigma) - \frac{1}{C'_3}\sqrt{\mathbb{E}\left[(x^Tv)^4\right]}.
\end{align*}
Since $x$ has bounded $4^{\text{th}}$ moments, we have $\mathbb{E}\left[(x^Tv)^4\right] \leq \widetilde{C}_4\mathbb{E}\left[(x^Tv)^2\right]^2\leq \widetilde{C}_4\lambda_{\max}(\Sigma)^2$. Hence, we obtain
\begin{align*}
v^T\left(\mathbb{E}[xx^T] - \mathbb{E}[xx^T\mathbbm{1}_{A^c}]\right)v \geq \lambda_{\min}(\Sigma) - \frac{\lambda_{\max}(\Sigma)\sqrt{\widetilde{C}_4}}{C'_3} \geq \frac{C'_1}{p} - \frac{C'_2\sqrt{\widetilde{C}_4}}{C'_3p} = \frac{C'_1}{2p}. 
\end{align*}
Since $||v||_2 = 1$ was arbitrary, and using $||\mathcal{C}||_2 \leq C''_1\sqrt{p}$ and $C'_3 = \frac{2C'_2\sqrt{\widetilde{C}_4}}{C'_1}$, Jensen's inequality then implies
\begin{align*}
\nabla^2\mathcal{R}(\theta) &\succeq \frac{q^3C'_1}{4p\left(q^2 + \left(C'_3||\mathcal{C}||_2\sqrt{\lambda_{\max}(\Sigma)} + \mathbb{E}[|w|]\right)^2\right)^{3/2}}I_p \\
%
& \succeq \frac{q^3C'_1}{4p\left(q^2 + \left(C'_3C''_1\sqrt{C'_2} + \mathbb{E}[|w|]\right)^2\right)^{3/2}}I_p\\
&= \frac{q^3(C'_1)^4}{4p\left((C'_1)^2q^2 + \left(2C''_1C'_2\sqrt{C'_2\widetilde{C}_4} + C'_1\mathbb{E}[|w|]\right)^2\right)^{3/2}}I_p\\
&\succeq \frac{q^3(C'_1)^4}{4p\left((C'_1)^2q^2 + 8(C''_1)^2(C'_2)^3\widetilde{C}_4 + 2(C'_1)^2\mathbb{E}[|w|]^2\right)^{3/2}}I_p\\
&\succeq \frac{q^3(C'_1)^4}{4p\left((C'_1)^2q^2 + 8(C''_1)^2(C'_2)^3\widetilde{C}_4 + 2(C'_1)^2\sigma_2^2\right)^{3/2}}I_p,
\end{align*}
as wanted.
%Since $\theta \in \mathcal{C}$ was arbitrary, we have the desired strong convexity over $\mathcal{C}$ as well.

Finally, for (4), since $\theta^* \in \mathcal{C}$, $x \indep w$, and $\mathbb{E}[x] = 0$, we have 
\begin{align*}
\nabla \mathcal{R}(\theta^*) = \mathbb{E}\left[\psi_q(x^T(\theta^* - \theta^*) + w)x\right] = \mathbb{E}[\psi_q(w)x] = \mathbb{E}[\psi_q(w)]\mathbb{E}[x] = 0,
\end{align*}
as required.
\end{proof}

\end{document}
